\lecture{25+}{January, 21 2026}{Exam 2024}


\begin{problemwithimage}{e2024_1}{1}
  For the static system shown in Figure 1, determine the gage pressure in \unit{\kPa} at point $a$.
  Liquid $A$ has an $\mathrm{SG} = \num{1,2} $ and liquid $B$ has an $\mathrm{SG} = \num{0,75}$.
  The liquid surrounding point $a$ is water ($\rho = \qty{1000}{\kg\per\m^3} $) and the tank on the left is open to atmosphere
\end{problemwithimage}

\begin{assumptions}
  \item Hydrostatic equilibrium
  \item All fluids are at rest and will not mix and will not exert any capillary effects on each other
  \item All liquids have constant densities 
  \item The left tank is open to the atmosphere, i.e. $p_b = 0$ gauge at the surface
\end{assumptions}

\begin{derivation}
  The heights of the different liquid-interfaces is understood as:
  \begin{itemize}
    \item Free surface of $B$ is at $z = \qty{0,9}{\m}$
    \item $AB$ interface is at $z = \qty{0,4}{\m}$
    \item $AW$ interface is at $z = \qty{0,4}{\m} + \qty{0,25}{\m} = \qty{0,65}{\m}$
    \item Point $a$ is at $z = \qty{0,125}{\m}$
  \end{itemize}
  We perform a ``hydrostatic walk'' from the free surface to point $a$.
  From the free surface to the $AB$ interface we move vertically downwards a distance if $B_z - AB = \qty{0,9}{\m} - \qty{0,4}{\m} = \qty{0,5}{\m}$. 
  This corresponds to a pressure change of:
  \begin{equation*}
    p_{AB} = \rho_B g \cdot (B_z - AB) = \num{0,75} \cdot \qty{1000}{\kg\per\m^3} \cdot \qty{9,82}{\m\per\s^2} \cdot \qty{0,5}{\m}   = \qty{3.6825}{\kilo\pascal}
  .\end{equation*}
  Then we ``move'' a distance of $AB - AW = \qty{0,4}{\m} - \qty{0,65}{\m} = -\qty{0,25}{\m}$ which corresponds to a pressure
  \begin{equation*}
    p_{AW} = p_{AB}- \rho_{A} \cdot g (AB - AW) = \qty{3,6825}{\kPa} - \num{1,2} \cdot \qty{1000}{\kg\per\m^3} \cdot \qty{9,82}{\m\per\s^2} \cdot \qty{0,25}{\m} = \qty{0.7365}{\kilo\pascal}
  .\end{equation*}
  Then we move a distance of $AW - a_z = \qty{0,65}{\m} - \qty{0,125}{\m} =  \qty{0,525}{\m}$ which corresponds to a pressure
  \begin{equation*}
    p_a = p_{AW} + \rho_w g (AW - a_z) = \qty{0,7365}{\kPa} + \qty{1000}{\kg\per\m^3} \cdot \qty{9,82}{\m\per\s^2} \cdot \qty{0,525}{\m}  = \qty{5.892}{\kilo\pascal}
  .\end{equation*}
  Hence, the gauge pressure at $a$ is about $p_a = \qty{5,89}{\kPa} (\text{gage})$.  
\end{derivation}

\finalresult{
p_a = \qty{5,89}{\kPa} 
}

\begin{problemwithimage}{e2024_2}{2}
  Water ($\rho = \qty{1000}{\kg\per\m^3}$) is flowing steadily through the \ang{180} elbow shown in Figure 2. 
  At the inlet to the elbow the gage pressure is \qty{103}{\kPa}. 
  The water discharges into atmospheric pressure. 
  Assume that the flow properties are uniform over the inlet and outlet areas; $A_1 = \qty{2500}{\mm^2}$, $A_2 = \qty{650}{\mm^2 }$, and $V_1 = \qty{3}{\m\per\s}$ are given.
  Find the horizontal component of force required to hold the elbow in place.
\end{problemwithimage}

\begin{assumptions}
  \item Steady, uniform flow at each section
  \item Incompressible fluid
  \item Negligible wall shear
\end{assumptions}

\begin{derivation}
  Due to mass conservation, the mass flow rate must be constant, $Q = \text{const.}$
  I.e.
  \begin{equation*}
    Q = A_1 V_1 = \qty{2500}{\mm^2} \cdot \qty{3}{\m\per\s} = \qty{0.0075}{\meter^3\per\second} \quad \implies \quad V_2 = \frac{Q}{A_2} = \frac{\qty{0,0075}{\m^3\per\second}}{\qty{650}{\mm^2} } = \qty{11.538}{\meter\per\second}
  .\end{equation*}
  As the outlet points left $V_{2x} = - \qty{11,54}{\m\per\s}$-
  This means the mass flow is
  \begin{equation*}
    \dot{m} = \rho Q = \qty{1000}{\kg\per\m^3} \cdot \qty{0,0075}{\m^3\per\s}  = \qty{7.5}{\kilogram\per\second}
  .\end{equation*}
  
  For steady flow, the linear momentum balance in $x$ is:
  \begin{equation*}
    \sum F_x = \int_{\mathrm{CS}} u \rho \Vec{V} \cdot \mathrm{d} \Vec{A} = \dot{m} \left( V_{2x} - V_{1x} \right) = \qty{7,5}{\kg\per\s} \cdot \left( - \qty{11,538}{\m\per\s} - \qty{3}{\m\per\s}  \right) = \qty{-109.035}{\newton}
  .\end{equation*}
  The negative sign indicates, that the net external force on the fluid points to the left.

  At the inlet the pressure acts normal into the elbow as the outward normal points to the left:
  \begin{equation*}
    F_{p_1, x} = p_1 A_1
  .\end{equation*}
  At the outlet the water discharges to the atmosphere so here the pressure $F_{p_2, x} = 0$ as $p_2 = 0$. This means, that the full force balance on the elbow in the $x$-direction is
  \begin{equation*}
    \sum F_x = p_1 A_1 + F_{\text{wall}, x}
  .\end{equation*}
  The sum of forces in the $x$-direction was previously found to be $\sum F_x = - \qty{109,035}{\N}$, meaning
  \begin{equation*}
    p_1 A_1 + F_{\text{wall}, x} = - \qty{109,035}{\N} \implies F_{\text{wall}, x} = - \qty{109,035}{\N} - p_1 A_1 = -\qty{109,035}{\N} - \qty{103}{\kPa} \cdot \qty{2500}{\mm^2}  = \qty{-0.36653500000000006}{\kilo\newton}
  .\end{equation*}
  Hence, the elbow will be pushed into the wall with a force of $F_{\text{fluid}} = \qty{366,5}{\N}$  and the wall must supply a force $F_{\text{wall}, x} = - \qty{366,5}{\N}$ 
\end{derivation}

\finalresult{
  F_{\text{wall}, x} = - \qty{366,5}{\N} 
}

\begin{problem}{3}
  Consider the velocity field $\Vec{V} = A \left( x^{4} - 6x^2 y^2 + y^{4} \right) \hat{i} + A \left( 4xy^3 - 4x^3 y \right) \hat{j}$, where $A = \qty{0,25}{\per\m^3\per\s}$.
  Show that the flow field is incompressible. Calculate the acceleration of a fluid particle at point $(x,y) = (2,1)$.
\end{problem}

\begin{assumptions}
  \item As no $t$ components are present in the velocity field it is implied to be steady
\end{assumptions}

\begin{derivation}
   For an incompressible flow the continuity equation simplifies to
   \begin{equation*}
     \nabla \cdot \Vec{V} = 0 \implies \frac{\partial u}{\partial x} + \frac{\partial v}{\partial y} = 0
   .\end{equation*}
   We compute the partial derivatives:
   \begin{align*}
     u &= A \left( x^{4} - 6x^2 y^2 + y^{4} \right) & \implies& & \frac{\partial u}{\partial x} &= A \left( 4x^3 - 12 xy^2 \right) \\
     v &= A \left( 4xy^3 - 4x^3 y \right) & \implies& & \frac{\partial v}{\partial y} &= A \left( 12 xy^2 - 4x^3 \right)
   .\end{align*}
   I.e.
   \begin{equation*}
     \frac{\partial u}{\partial x} + \frac{\partial v}{\partial y} = A \left( 4x^3 - 4x^3 + 12xy^2 - 12xy^2 \right) = 0
   \end{equation*}
   meaning the flow field is incompressible everywhere.
   
   For the acceleration we use that:
   \begin{equation*}
     \Vec{a} = (\Vec{V} \cdot \nabla ) \Vec{V}
   \end{equation*}
   which, in Cartesian coordinates is
   \begin{equation*}
     a_x = u \frac{\partial u}{\partial x} + v \frac{\partial u}{\partial y}, \qquad a_y = u \frac{\partial v}{\partial x} + v \frac{\partial v}{\partial y} 
   .\end{equation*}
   We compute the two missing derivatives as:
   \begin{align*}
     \frac{\partial u}{\partial y} &= A \left( 4 y^3 - 12 x^2 y \right) \\
     \frac{\partial v}{\partial x} &= A \left( 4y^3 - 12x^2 y \right)
   .\end{align*}
   So the acceleration components are
   \begin{align*}
     a_x &= u \frac{\partial u}{\partial x} + v \frac{\partial u}{\partial y}  \\
     &= A^2 \left( x^{4} - 6 x^2 y^2 + y^{4} \right) \left( 4 x^3 - 12 x y^2 \right) + A^2 \left( 4 x y^3 - 4 x^3 y \right) \left( 4 y^3 - 12 x^2 y \right) \\
     &=4 A^2 x \left(x^2+y^2\right)^3\\
     a_y &= u \frac{\partial v}{\partial x} + v \frac{\partial v}{\partial y}  \\
     &= A^2 \left( x^{4} - 6 x^2 y^2 + y^{4} \right) \left( 4 y^3 - 12 x^2 y \right) + A^2 \left( 4 x y^3 - 4 x^3 y \right) \left( 12 x y^2 - 4 x^3 \right) \\
     &= 4 A^2 y \left(x^2+y^2\right)^3
   .\end{align*}
   So,
   \begin{align*}
     a_x(2,1) &= 4 \cdot \left( \num{0,25}  \right)^2 \cdot 2 \left( 4 + 1  \right)^3 = 62.5 \\
     a_y(2,1) &= 4 \cdot \left( \num{0,25} \right)^2 \cdot 1\cdot  \left( 4 + 1 \right)^3    = 31.25
   .\end{align*}
   So,
   \begin{equation*}
     a(2,1) = \qty{62,5}{\m\per\s^2} \hat{i} + \qty{31,25}{\m\per\s^2} \hat{j}
   .\end{equation*}
\end{derivation}

\finalresult{
\begin{gathered}
  \frac{\partial u}{\partial x} + \frac{\partial v}{\partial y} = 0 \\
  a(2,1) = \qty{62,5}{\m\per\s^2} \hat{i} + \qty{31,25}{\m\per\s^2} \hat{j}
.\end{gathered}
}

\begin{problem}{4}
  A model airplane has a span of \qty{120}{\cm}  and a chord of \qty{25}{\cm}.
  The model flies in air at standard atmospheric conditions (i.e., $\rho = \qty{1,2}{\kg\per\m^3}, \nu = \qty{1,5e-5}{\m^2\per\s} $).
  The wing can be regarded as a flat plate so far as drag is concerned.
  At what speed will a transition to a turbulent boundary layer start?
  What is the total drag force on the wing just before the transition to turbulence happens?
\end{problem}

\begin{assumptions}
  \item The wing is flat and has zero thickness
  \item Flow is parallel to the wing, meaning to lift effects are present
  \item The freestream is uniform incompressible ($M \ll \num{0,3}$), and low-turbulence
  \item The surface of the wing is assumed smooth
  \item Transition to turbulent flow happens at $\mathrm{Re}_{x, \mathrm{cr}} = \num{5e5}$
\end{assumptions}

\begin{derivation}
  The Reynolds number for the flow over the flat plate (the wing) is
  \begin{equation*}
    \mathrm{Re}_{x} = \frac{V x}{\nu}, \qquad \mathrm{Re}_{x, \mathrm{cr}} \approx \num{5e5}
  .\end{equation*}
  Due to the length-parameter of the Reynolds number transition to turbulent flow will happen at the trailing edge first.
  Inserting the chord length as our characteristic length and solving for the critical speed we obtain:
  \begin{equation*}
    V_{\mathrm{cr}} = \frac{\mathrm{Re}_{x, \mathrm{cr}} \nu}{L} = \frac{\num{5e5} \cdot \qty{1,5e-5}{\m^2\per\s}}{\qty{0,25}{\m} } = \qty{30}{\m\per\s} 
  .\end{equation*}
  
  The drag coefficient for flow with free stream velocity $V$, over a flat plate of length $L$ and width $b$ is:
  \begin{equation*}
    C_d = \frac{\num{1,33} }{\sqrt{\mathrm{Re}_L}} = \frac{\num{1,33} }{\sqrt{\num{5e5} }} = \num{1,88e-3} 
  .\end{equation*}
  
  The drag force $F_D$ is
  \begin{equation*}
    F_D = \frac{1}{2} \rho V^2 A C_D
  .\end{equation*}
  The area which the friction applies over $A$ is in this case both the bottom and top of the wing, hence:
  \begin{equation*}
    A = 2 \cdot \qty{120}{\cm} \cdot \qty{25}{\cm}  = \qty{0,6}{\meter^2}
  .\end{equation*}
  So
  \begin{equation*}
    F_D = \frac{1}{2} \cdot \qty{1,2}{\kg\per\m^3} \cdot \left( \qty{30}{\m\per\s}  \right)^2 \cdot \qty{0,6}{\m^2} \cdot \num{1,88e-3}  = \qty{0,61}{\N} 
  .\end{equation*}
\end{derivation}

\finalresult{
\begin{aligned}
  V_{\mathrm{cr}} &= \qty{30}{\m\per\s} \\
  F_D &= \qty{0,61}{\N} 
.\end{aligned}
}
